\documentclass[a4paper,12pt]{report}

% ── Packages ──────────────────────────────────────────────────
\usepackage[utf8]{inputenc}
\usepackage[T1]{fontenc}
\usepackage[french]{babel}
\usepackage{geometry}
\geometry{margin=2.5cm}
\usepackage{graphicx}
\usepackage{hyperref}
\usepackage{listings}
\usepackage{xcolor}
\usepackage{fancyhdr}
\usepackage{titlesec}
\usepackage{enumitem}
\usepackage{tabularx}
\usepackage{booktabs}
\usepackage{longtable}
\usepackage{tcolorbox}
\usepackage{lmodern}

% ── Couleurs ──────────────────────────────────────────────────
\definecolor{primary}{HTML}{1E40AF}
\definecolor{codebg}{HTML}{F8F9FA}
\definecolor{codeframe}{HTML}{DEE2E6}
\definecolor{keyword}{HTML}{D73A49}
\definecolor{string}{HTML}{032F62}
\definecolor{comment}{HTML}{6A737D}
\definecolor{foldercolor}{HTML}{E65100}
\definecolor{filecolor}{HTML}{1565C0}

% ── Configuration listings ────────────────────────────────────
\lstset{
  backgroundcolor=\color{codebg},
  frame=single,
  rulecolor=\color{codeframe},
  basicstyle=\ttfamily\small,
  keywordstyle=\color{keyword}\bfseries,
  stringstyle=\color{string},
  commentstyle=\color{comment}\itshape,
  breaklines=true,
  showstringspaces=false,
  tabsize=2,
  numbers=left,
  numberstyle=\tiny\color{gray},
  numbersep=8pt,
  xleftmargin=15pt,
  framexleftmargin=15pt,
}

% ── En-tête / Pied de page ───────────────────────────────────
\pagestyle{fancy}
\fancyhf{}
\fancyhead[L]{\small\textcolor{primary}{SGRH Intelligent}}
\fancyhead[R]{\small\textcolor{primary}{Documentation Technique}}
\fancyfoot[C]{\thepage}
\renewcommand{\headrulewidth}{0.5pt}

% ── Titre sections colorés ───────────────────────────────────
\titleformat{\chapter}[display]
  {\normalfont\huge\bfseries\color{primary}}
  {\chaptertitlename\ \thechapter}{20pt}{\Huge}
\titleformat{\section}
  {\normalfont\Large\bfseries\color{primary}}
  {\thesection}{1em}{}
\titleformat{\subsection}
  {\normalfont\large\bfseries\color{primary!80!black}}
  {\thesubsection}{1em}{}
\titleformat{\subsubsection}
  {\normalfont\normalsize\bfseries\color{primary!60!black}}
  {\thesubsubsection}{1em}{}

% ── Boîtes colorées ──────────────────────────────────────────
\newtcolorbox{infobox}[1][]{
  colback=blue!5!white,
  colframe=primary,
  fonttitle=\bfseries,
  title=#1,
  rounded corners,
  boxrule=0.5pt,
}

\newtcolorbox{warningbox}[1][]{
  colback=orange!5!white,
  colframe=orange!80!black,
  fonttitle=\bfseries,
  title=#1,
  rounded corners,
  boxrule=0.5pt,
}

\newtcolorbox{folderbox}{
  colback=orange!3!white,
  colframe=foldercolor,
  rounded corners,
  boxrule=0.5pt,
  left=5pt, right=5pt, top=3pt, bottom=3pt,
}

\newtcolorbox{filebox}{
  colback=blue!3!white,
  colframe=filecolor,
  rounded corners,
  boxrule=0.5pt,
  left=5pt, right=5pt, top=3pt, bottom=3pt,
}

% ══════════════════════════════════════════════════════════════
\begin{document}

% ── Page de garde ─────────────────────────────────────────────
\begin{titlepage}
\centering
\vspace*{2cm}

{\Large Université Hassan II de Casablanca}\\[0.3cm]
{\large Faculté des Sciences Ben M'Sick}\\[0.3cm]
{\large Département Mathématiques et Informatique}\\[2cm]

\rule{\textwidth}{1.5pt}\\[0.5cm]
{\Huge\bfseries\textcolor{primary}{Documentation Technique}}\\[0.3cm]
{\LARGE\textcolor{primary}{SGRH Intelligent}}\\[0.2cm]
{\large Système de Gestion des Ressources Humaines\\avec Détection d'Anomalies par IA et Analyse OCR}\\[0.3cm]
\rule{\textwidth}{1.5pt}\\[2cm]

{\large\textbf{PFE -- Licence 3\textsuperscript{ème} Année Développement Informatique}}\\[0.3cm]
{\large Année Universitaire 2025 -- 2026}\\[2cm]

\begin{tabular}{ll}
\textbf{Encadré par :} & Pr. Ichrak BENAMRI \\[0.3cm]
\textbf{Réalisé par :} & MANSOUR Aymane \\
                        & FOUKOU Issa \\
                        & BADIA Salahedine \\
\end{tabular}

\vfill
{\small Faculté des Sciences Ben M'Sick --- Casablanca}
\end{titlepage}

% ── Table des matières ────────────────────────────────────────
\tableofcontents
\newpage

% ══════════════════════════════════════════════════════════════
\chapter{Vue d'ensemble du projet}
% ══════════════════════════════════════════════════════════════

\section{Présentation générale}

Le \textbf{SGRH Intelligent} est une application web complète de gestion des ressources humaines. Elle se distingue des SGRH classiques par l'intégration de deux modules innovants :

\begin{itemize}[leftmargin=2cm]
  \item[\textbf{IA}] Détection automatique d'anomalies comportementales (retards, absences suspectes) via les algorithmes \textbf{Isolation Forest} et \textbf{K-Means}.
  \item[\textbf{OCR}] Extraction automatique de texte depuis des documents administratifs numérisés (certificats médicaux, contrats) via le moteur \textbf{Tesseract}.
\end{itemize}

\section{Architecture globale}

L'application suit une \textbf{architecture trois tiers enrichie d'un microservice IA} :

\begin{center}
\begin{tabularx}{\textwidth}{|l|l|l|X|}
\hline
\textbf{Couche} & \textbf{Technologie} & \textbf{Port} & \textbf{Rôle} \\
\hline
Présentation & React JS + Vite & 3000 & Interface utilisateur, tableau de bord, formulaires \\
\hline
Backend & Spring Boot / Java 21 & 8080 & API REST, logique métier, sécurité JWT/RBAC \\
\hline
Microservice IA & FastAPI / Python & 8000 & Isolation Forest, K-Means, génération de rapports \\
\hline
OCR & Tesseract 5.x & -- & Extraction de texte (PDF, images) \\
\hline
Données & MySQL 8 & 3306 & Base de données relationnelle \\
\hline
\end{tabularx}
\end{center}

\begin{infobox}[Flux de communication]
\begin{enumerate}
  \item Le \textbf{Frontend} envoie des requêtes HTTP REST au \textbf{Backend} avec un token JWT dans le header \texttt{Authorization}.
  \item Le \textbf{Backend} traite la logique métier et accède à \textbf{MySQL} via JPA/Hibernate.
  \item Pour l'analyse IA, le Backend appelle le \textbf{Microservice FastAPI} via HTTP POST.
  \item Pour l'OCR, le Backend invoque \textbf{Tesseract} localement sur les fichiers uploadés.
  \item Les résultats sont retournés au Frontend pour affichage dans le tableau de bord.
\end{enumerate}
\end{infobox}

\section{Arborescence complète du projet}

\begin{lstlisting}[language={},title={\textbf{Structure des dossiers}},morekeywords={frontend,backend,microservice-ia,database,documentation,src,main,java,ma,fsb,sgrh,components,context,pages,services,entity,repository,controller,security,service,dto,resources}]
PFE_Slah/
|-- frontend/                    # Interface utilisateur React
|   |-- index.html               # Point d'entree HTML
|   |-- package.json             # Dependances NPM
|   |-- vite.config.js           # Configuration Vite
|   |-- tailwind.config.js       # Configuration TailwindCSS
|   |-- postcss.config.js        # Configuration PostCSS
|   +-- src/
|       |-- main.jsx             # Demarrage de React
|       |-- App.jsx              # Routeur principal
|       |-- index.css            # Styles TailwindCSS
|       |-- context/
|       |   +-- AuthContext.jsx  # Gestion authentification
|       |-- services/
|       |   +-- api.js           # Client Axios + intercepteurs
|       |-- components/
|       |   |-- Layout.jsx       # Mise en page generale
|       |   |-- Sidebar.jsx      # Menu lateral
|       |   +-- Header.jsx       # Barre superieure
|       +-- pages/
|           |-- LoginPage.jsx    # Page de connexion
|           |-- DashboardPage.jsx # Tableau de bord
|           |-- EmployesPage.jsx # Gestion des employes
|           |-- CongesPage.jsx   # Gestion des conges
|           |-- PresencesPage.jsx # Suivi des presences
|           |-- PaiePage.jsx     # Fiches de paie
|           |-- DocumentsPage.jsx # Documents & OCR
|           |-- IAPage.jsx       # Analyse IA
|           +-- UsersPage.jsx    # Gestion utilisateurs
|-- backend/                     # API REST Spring Boot
|   |-- pom.xml                  # Dependances Maven
|   +-- src/main/
|       |-- resources/
|       |   +-- application.yml  # Configuration Spring
|       +-- java/ma/fsb/sgrh/
|           |-- SgrhApplication.java # Classe principale
|           |-- entity/          # Entites JPA (7 classes)
|           |-- repository/      # Repositories Spring Data
|           |-- controller/      # Controleurs REST (6)
|           |-- service/         # Services metier
|           |-- security/        # JWT + Spring Security
|           +-- dto/             # Objets de transfert
|-- microservice-ia/             # Module IA Python
|   |-- main.py                  # Serveur FastAPI
|   +-- requirements.txt        # Dependances pip
|-- database/                    # Scripts SQL
|   |-- schema.sql               # Schema de la BDD
|   +-- seed.sql                 # Donnees de test
+-- README.md                   # Guide de demarrage
\end{lstlisting}


% ══════════════════════════════════════════════════════════════
\chapter{Base de données (\texttt{database/})}
% ══════════════════════════════════════════════════════════════

\section{\texttt{schema.sql} --- Schéma de la base de données}

Ce fichier crée la base \texttt{sgrh\_intelligent} et définit les 7 tables principales, conformément au diagramme de classes UML et à la méthode Merise (MCD/MLD).

\subsection{Tables créées}

\begin{center}
\begin{tabularx}{\textwidth}{|l|X|}
\hline
\textbf{Table} & \textbf{Description} \\
\hline
\texttt{utilisateur} & Comptes utilisateurs avec email, mot de passe hashé (BCrypt), rôle (\texttt{ADMIN}, \texttt{RH}, \texttt{EMPLOYE}) et statut actif/inactif. \\
\hline
\texttt{employe} & Informations personnelles et professionnelles : matricule, nom, prénom, poste, département, salaire, solde de congé. Liée à \texttt{utilisateur} via clé étrangère (relation 1:1). \\
\hline
\texttt{conge} & Demandes de congé avec dates, type (\texttt{ANNUEL}, \texttt{MALADIE}, \texttt{SANS\_SOLDE}), statut (\texttt{EN\_ATTENTE}, \texttt{VALIDE}, \texttt{REFUSE}) et motif. \\
\hline
\texttt{presence} & Pointages quotidiens : date, heure d'arrivée, heure de départ, indicateur de retard. Contrainte d'unicité sur (employé, date). \\
\hline
\texttt{fiche\_paie} & Fiches de paie mensuelles avec salaire brut et net. Contrainte d'unicité sur (employé, mois, année). \\
\hline
\texttt{document} & Documents uploadés : nom de fichier, type, chemin physique, texte extrait par OCR, statut de traitement. \\
\hline
\texttt{anomalie} & Résultats de l'analyse IA : score d'isolation, type d'anomalie, niveau de risque (\texttt{FAIBLE}, \texttt{MOYEN}, \texttt{ELEVE}), label de cluster. \\
\hline
\end{tabularx}
\end{center}

\subsection{Index de performance}
Quatre index sont créés pour accélérer les requêtes fréquentes :
\begin{itemize}
  \item \texttt{idx\_conge\_employe\_statut} : recherche rapide des congés par employé et statut.
  \item \texttt{idx\_presence\_date} : filtrage des présences par date.
  \item \texttt{idx\_anomalie\_risque} : tri des anomalies par niveau de risque et date.
  \item \texttt{idx\_employe\_departement} : filtrage par département.
\end{itemize}

\subsection{Conformité loi 09-08}
La base est configurée en \texttt{utf8mb4} pour supporter le français et l'arabe. Les mots de passe sont stockés en hash BCrypt (jamais en clair). Les identifiants personnels peuvent être anonymisés avant tout traitement IA.

\section{\texttt{seed.sql} --- Données de test}

Ce fichier insère des données réalistes pour tester l'application :
\begin{itemize}
  \item \textbf{7 utilisateurs} : 1 admin, 1 RH, 5 employés (mot de passe BCrypt pour \texttt{password123}).
  \item \textbf{6 employés} avec des postes variés (développeur, designer, comptable, marketing, technicien).
  \item \textbf{5 congés} couvrant les trois statuts (validé, en attente, refusé).
  \item \textbf{12 présences} sur une semaine type, incluant des retards.
  \item \textbf{5 fiches de paie} pour mai 2025.
\end{itemize}


% ══════════════════════════════════════════════════════════════
\chapter{Frontend (\texttt{frontend/})}
% ══════════════════════════════════════════════════════════════

\section{Fichiers de configuration}

\subsection{\texttt{package.json}}
Fichier de configuration NPM qui déclare :
\begin{itemize}
  \item Le nom du projet (\texttt{sgrh-intelligent-frontend}) et sa version.
  \item Les \textbf{dépendances de production} : \texttt{react}, \texttt{react-dom}, \texttt{react-router-dom} (routage SPA), \texttt{axios} (appels HTTP), \texttt{recharts} (graphiques), \texttt{lucide-react} (icônes), \texttt{date-fns} (dates), \texttt{react-hot-toast} (notifications).
  \item Les \textbf{dépendances de développement} : \texttt{vite} (bundler ultra-rapide), \texttt{tailwindcss} (framework CSS utilitaire), \texttt{autoprefixer}, \texttt{postcss}.
  \item Les \textbf{scripts} : \texttt{dev} (serveur de développement), \texttt{build} (production), \texttt{preview}.
\end{itemize}

\subsection{\texttt{vite.config.js}}
Configuration du bundler Vite :
\begin{itemize}
  \item Active le plugin \texttt{@vitejs/plugin-react} pour le support JSX.
  \item Définit un alias \texttt{@} pointant vers \texttt{./src} pour des imports propres.
  \item Configure un \textbf{proxy} : toutes les requêtes \texttt{/api} sont redirigées vers \texttt{http://localhost:8080} (backend Spring Boot), évitant les problèmes CORS en développement.
  \item Le serveur de développement écoute sur le \textbf{port 3000}.
\end{itemize}

\subsection{\texttt{tailwind.config.js}}
Configuration de TailwindCSS :
\begin{itemize}
  \item \texttt{content} : scanne tous les fichiers \texttt{.jsx/.tsx} dans \texttt{src/} pour générer uniquement les classes CSS utilisées (purge automatique).
  \item \texttt{colors.primary} : palette de bleus personnalisée (50 à 900) pour l'identité visuelle.
  \item \texttt{colors.sidebar} : couleurs sombres (\texttt{\#0f172a}) pour le menu latéral.
  \item \texttt{fontFamily} : police \texttt{Inter} pour une typographie moderne.
\end{itemize}

\subsection{\texttt{postcss.config.js}}
Active les plugins PostCSS nécessaires : \texttt{tailwindcss} (compilation des classes utilitaires) et \texttt{autoprefixer} (ajout automatique des préfixes navigateurs).

\subsection{\texttt{index.html}}
Point d'entrée HTML de l'application :
\begin{itemize}
  \item Charge la police \textbf{Inter} depuis Google Fonts.
  \item Contient un \texttt{<div id="root">} où React monte l'application.
  \item Référence \texttt{/src/main.jsx} comme module d'entrée JavaScript.
\end{itemize}

\section{Fichiers source (\texttt{src/})}

\subsection{\texttt{main.jsx} --- Point d'entrée React}
Ce fichier initialise l'application React :
\begin{enumerate}
  \item Crée la racine React avec \texttt{ReactDOM.createRoot}.
  \item Enveloppe l'application dans \texttt{BrowserRouter} (routage côté client).
  \item Enveloppe dans \texttt{AuthProvider} (contexte d'authentification global).
  \item Ajoute le composant \texttt{Toaster} pour les notifications toast.
  \item Active \texttt{StrictMode} pour détecter les problèmes potentiels.
\end{enumerate}

\subsection{\texttt{App.jsx} --- Routeur principal}
Définit toutes les routes de l'application avec \textbf{React Router v6} :

\begin{center}
\begin{tabularx}{\textwidth}{|l|l|X|}
\hline
\textbf{Route} & \textbf{Composant} & \textbf{Accès} \\
\hline
\texttt{/login} & \texttt{LoginPage} & Public (redirige si déjà connecté) \\
\hline
\texttt{/dashboard} & \texttt{DashboardPage} & Tous les rôles authentifiés \\
\hline
\texttt{/employes} & \texttt{EmployesPage} & ADMIN, RH uniquement \\
\hline
\texttt{/conges} & \texttt{CongesPage} & Tous les rôles authentifiés \\
\hline
\texttt{/presences} & \texttt{PresencesPage} & Tous les rôles authentifiés \\
\hline
\texttt{/paie} & \texttt{PaiePage} & ADMIN, RH uniquement \\
\hline
\texttt{/documents} & \texttt{DocumentsPage} & Tous les rôles authentifiés \\
\hline
\texttt{/ia} & \texttt{IAPage} & ADMIN, RH uniquement \\
\hline
\texttt{/utilisateurs} & \texttt{UsersPage} & ADMIN uniquement \\
\hline
\end{tabularx}
\end{center}

Le composant \texttt{ProtectedRoute} vérifie :
\begin{enumerate}
  \item Si l'utilisateur est authentifié (sinon, redirection vers \texttt{/login}).
  \item Si l'utilisateur a le rôle requis (sinon, redirection vers \texttt{/dashboard}).
\end{enumerate}

\subsection{\texttt{index.css} --- Styles TailwindCSS}
Fichier CSS principal qui :
\begin{itemize}
  \item Importe les trois couches Tailwind : \texttt{@tailwind base}, \texttt{components}, \texttt{utilities}.
  \item Définit des \textbf{composants réutilisables} dans \texttt{@layer components} :
  \begin{itemize}
    \item \texttt{.btn-primary}, \texttt{.btn-secondary}, \texttt{.btn-danger} : boutons stylisés.
    \item \texttt{.card} : conteneur blanc avec ombre et coins arrondis.
    \item \texttt{.input-field} : champs de formulaire avec focus ring bleu.
    \item \texttt{.badge-success}, \texttt{.badge-warning}, \texttt{.badge-danger}, \texttt{.badge-info} : étiquettes colorées.
  \end{itemize}
\end{itemize}

\section{Contexte (\texttt{context/})}

\subsection{\texttt{AuthContext.jsx} --- Gestion de l'authentification}
Ce fichier implémente le \textbf{Context API} de React pour gérer l'état d'authentification globalement :

\begin{itemize}
  \item \textbf{\texttt{user}} : objet contenant \texttt{id}, \texttt{email}, \texttt{role} de l'utilisateur connecté.
  \item \textbf{\texttt{login(userData, token)}} : sauvegarde l'utilisateur et le token JWT dans \texttt{localStorage}.
  \item \textbf{\texttt{logout()}} : supprime les données de \texttt{localStorage} et remet \texttt{user} à \texttt{null}.
  \item \textbf{\texttt{getToken()}} : retourne le token JWT stocké.
  \item \textbf{\texttt{loading}} : état booléen qui évite un flash de la page login pendant la vérification initiale.
\end{itemize}

Au montage du composant (\texttt{useEffect}), le contexte vérifie si un utilisateur et un token sont déjà stockés dans \texttt{localStorage} pour maintenir la session.

\section{Services (\texttt{services/})}

\subsection{\texttt{api.js} --- Client HTTP Axios}
Ce fichier centralise \textbf{tous les appels API} de l'application :

\begin{itemize}
  \item Crée une instance Axios avec l'URL de base \texttt{http://localhost:8080/api}.
  \item \textbf{Intercepteur de requête} : ajoute automatiquement le header \texttt{Authorization: Bearer <token>} à chaque requête si un token JWT existe.
  \item \textbf{Intercepteur de réponse} : si le serveur retourne un code 401 (non autorisé), déconnecte l'utilisateur et redirige vers \texttt{/login}.
  \item Exporte 7 \textbf{services} organisés par module :
  \begin{itemize}
    \item \texttt{authService} : \texttt{login()}, \texttt{register()}, \texttt{me()}.
    \item \texttt{employeService} : CRUD complet (\texttt{getAll}, \texttt{getById}, \texttt{create}, \texttt{update}, \texttt{delete}).
    \item \texttt{congeService} : \texttt{getAll}, \texttt{getByEmploye}, \texttt{create}, \texttt{valider}, \texttt{refuser}.
    \item \texttt{presenceService} : \texttt{getAll}, \texttt{getByEmploye}, \texttt{pointer}.
    \item \texttt{paieService} : \texttt{getAll}, \texttt{getByEmploye}, \texttt{generer}.
    \item \texttt{documentService} : \texttt{getAll}, \texttt{upload} (multipart/form-data), \texttt{getById}.
    \item \texttt{iaService} : \texttt{detecterAnomalies}, \texttt{segmenter}, \texttt{getRapport}, \texttt{getAnomalies}.
  \end{itemize}
\end{itemize}

\section{Composants (\texttt{components/})}

\subsection{\texttt{Layout.jsx} --- Mise en page}
Composant de mise en page principal qui structure l'interface en deux zones :
\begin{itemize}
  \item \textbf{Sidebar} (à gauche, largeur fixe 256px) : menu de navigation.
  \item \textbf{Zone principale} (à droite, flexible) : contient le \texttt{Header} en haut et le contenu de la page (\texttt{<Outlet />} de React Router) en dessous.
\end{itemize}
Utilise \texttt{flexbox} avec \texttt{h-screen overflow-hidden} pour occuper tout l'écran.

\subsection{\texttt{Sidebar.jsx} --- Menu latéral}
Menu de navigation vertical avec fond sombre (\texttt{\#0f172a}) :
\begin{itemize}
  \item \textbf{Logo} : icône Brain avec le texte « SGRH Intelligent ».
  \item \textbf{Navigation} : liste de liens filtrés selon le rôle de l'utilisateur. Chaque lien a une icône Lucide et un label. Le lien actif est surligné en bleu (\texttt{NavLink} avec \texttt{isActive}).
  \item \textbf{Pied} : avatar avec l'initiale de l'email, rôle affiché, et bouton de déconnexion.
  \item \textbf{Filtrage par rôle} : chaque élément de navigation déclare les rôles autorisés. Seuls les éléments correspondant au rôle de l'utilisateur s'affichent.
\end{itemize}

\subsection{\texttt{Header.jsx} --- Barre supérieure}
Barre horizontale contenant :
\begin{itemize}
  \item Un champ de \textbf{recherche} avec icône loupe.
  \item Une \textbf{cloche de notifications} avec indicateur rouge.
  \item L'\textbf{avatar} de l'utilisateur avec son nom et rôle traduit en français.
\end{itemize}

\section{Pages (\texttt{pages/})}

\subsection{\texttt{LoginPage.jsx} --- Page de connexion}
Page divisée en deux panneaux :
\begin{itemize}
  \item \textbf{Panneau gauche} (visible sur écrans larges) : branding avec gradient bleu, logo, description du projet, et trois cartes « IA / OCR / RH ».
  \item \textbf{Panneau droit} : formulaire de connexion avec email et mot de passe, icônes, bouton de visibilité du mot de passe, et indicateur de chargement.
  \item \textbf{Mode démo} : trois boutons (Admin / RH / Employé) permettent de se connecter \textbf{sans backend} avec des données fictives. Idéal pour la soutenance.
\end{itemize}

\subsection{\texttt{DashboardPage.jsx} --- Tableau de bord}
Page d'accueil après connexion, composée de :
\begin{itemize}
  \item \textbf{4 cartes statistiques} : Total employés, Congés en attente, Présents aujourd'hui, Anomalies IA. Chaque carte affiche une tendance (hausse/baisse).
  \item \textbf{Graphique à barres} (Recharts) : congés par mois, ventilés par type (annuels, maladie, sans solde).
  \item \textbf{Graphique en anneau} : répartition des employés par département.
  \item \textbf{Graphique d'aire} : présences et retards de la semaine.
  \item \textbf{Liste d'activité récente} : 5 dernières actions du système avec horodatage.
\end{itemize}

\subsection{\texttt{EmployesPage.jsx} --- Gestion des employés}
Interface CRUD complète pour les employés :
\begin{itemize}
  \item \textbf{Barre de recherche} filtrant en temps réel par nom, matricule ou département.
  \item \textbf{Tableau} avec colonnes : avatar/nom, matricule (police monospace), poste, département (badge bleu), salaire (formaté MAD), statut (badge vert/rouge).
  \item \textbf{Actions} par ligne : voir, modifier, supprimer (icônes avec hover coloré).
  \item \textbf{Modal d'ajout} : formulaire en deux colonnes (nom, prénom, email, poste, département, salaire).
  \item \textbf{Bouton export} pour télécharger la liste.
\end{itemize}

\subsection{\texttt{CongesPage.jsx} --- Gestion des congés}
\begin{itemize}
  \item \textbf{Filtres par statut} : boutons « Tous / En attente / Validé / Refusé ».
  \item \textbf{Cartes de congé} : chaque demande affiche l'employé, les dates, le nombre de jours, le type, le statut (badge coloré) et le motif.
  \item \textbf{Actions RH} : boutons Valider (vert) et Refuser (rouge) visibles uniquement pour les rôles RH/Admin sur les congés en attente.
  \item \textbf{Modal de nouvelle demande} : formulaire avec sélecteurs de dates, type de congé et motif.
\end{itemize}

\subsection{\texttt{PresencesPage.jsx} --- Suivi des présences}
\begin{itemize}
  \item \textbf{Sélecteur de date} pour filtrer par jour.
  \item \textbf{4 cartes} : Présents, Absents, Retards, Taux de présence (\%).
  \item \textbf{Tableau des pointages} : employé, date, heure d'arrivée/départ (monospace), durée calculée, statut (« À l'heure » ou « En retard »).
\end{itemize}

\subsection{\texttt{PaiePage.jsx} --- Gestion de la paie}
\begin{itemize}
  \item \textbf{3 cartes} : masse salariale nette, salaire moyen, nombre de fiches générées.
  \item \textbf{Graphique à barres} : évolution de la masse salariale sur 6 mois (couleur violet).
  \item \textbf{Tableau des fiches} : employé, période, salaire brut, salaire net (en gras), bouton de téléchargement.
  \item \textbf{Bouton « Générer les fiches »} et \textbf{« Exporter PDF »}.
\end{itemize}

\subsection{\texttt{DocumentsPage.jsx} --- Documents \& OCR}
\begin{itemize}
  \item \textbf{3 cartes statistiques} : total documents, OCR réussi, en traitement.
  \item \textbf{Grille de documents} (2 colonnes) : chaque carte affiche le nom du fichier, le type, l'employé, le statut (Traité/En cours/Erreur), et un aperçu du texte extrait.
  \item \textbf{Modal d'upload} : zone de drag-and-drop avec sélecteur de type de document.
  \item \textbf{Modal de détail} : affiche toutes les métadonnées et le texte extrait complet par OCR.
\end{itemize}

\subsection{\texttt{IAPage.jsx} --- Analyse IA}
Page la plus riche visuellement :
\begin{itemize}
  \item \textbf{Bouton « Lancer l'analyse »} avec animation de chargement (spinner).
  \item \textbf{4 cartes} : anomalies détectées, taux de contamination (5\%), nombre de clusters (3), précision (92\%).
  \item \textbf{Scatter plot} (nuage de points) : visualisation des clusters K-Means avec couleurs (vert=Normal, jaune=Retardataire, rouge=Absent).
  \item \textbf{Graphique en anneau} : distribution des profils comportementaux.
  \item \textbf{Tableau des anomalies} : employé, barre de score visuelle, type d'anomalie, niveau de risque (badge), cluster assigné (pastille colorée), date.
\end{itemize}

\subsection{\texttt{UsersPage.jsx} --- Gestion des utilisateurs (Admin)}
\begin{itemize}
  \item \textbf{3 cartes} : nombre d'administrateurs, comptes actifs, comptes désactivés.
  \item \textbf{Tableau} : email, rôle (badge coloré par rôle), statut, dernière connexion.
  \item \textbf{Modal de création} : email, mot de passe, sélection du rôle.
\end{itemize}


% ══════════════════════════════════════════════════════════════
\chapter{Backend (\texttt{backend/})}
% ══════════════════════════════════════════════════════════════

\section{Configuration}

\subsection{\texttt{pom.xml} --- Dépendances Maven}
Le fichier POM déclare les dépendances du projet Spring Boot :

\begin{center}
\begin{tabularx}{\textwidth}{|l|X|}
\hline
\textbf{Dépendance} & \textbf{Rôle} \\
\hline
\texttt{spring-boot-starter-web} & Framework web, serveur Tomcat intégré, API REST \\
\hline
\texttt{spring-boot-starter-data-jpa} & ORM Hibernate, repositories auto-générés \\
\hline
\texttt{spring-boot-starter-security} & Spring Security, filtres d'authentification \\
\hline
\texttt{spring-boot-starter-validation} & Validation des DTOs avec annotations Jakarta \\
\hline
\texttt{mysql-connector-j} & Driver JDBC pour MySQL \\
\hline
\texttt{jjwt-api/impl/jackson} & Bibliothèque JWT pour génération et validation des tokens \\
\hline
\texttt{lombok} & Génération automatique de getters, setters, constructeurs \\
\hline
\texttt{spring-boot-starter-webflux} & WebClient pour appeler le microservice IA (HTTP réactif) \\
\hline
\end{tabularx}
\end{center}

\subsection{\texttt{application.yml} --- Configuration Spring}
Fichier YAML qui configure :
\begin{itemize}
  \item \textbf{DataSource} : URL MySQL (\texttt{jdbc:mysql://localhost:3306/sgrh\_intelligent}), identifiants.
  \item \textbf{JPA/Hibernate} : mode \texttt{ddl-auto: update} (création/modification automatique des tables).
  \item \textbf{JWT} : clé secrète de signature (256 bits), durée d'expiration (24h = 86400000ms).
  \item \textbf{Microservice IA} : URL du service FastAPI (\texttt{http://localhost:8000}).
  \item \textbf{OCR} : chemin vers l'exécutable Tesseract et langues (\texttt{fra+ara}).
  \item \textbf{Upload} : taille maximale des fichiers (10 Mo).
\end{itemize}

\subsection{\texttt{SgrhApplication.java} --- Classe principale}
Point d'entrée de l'application Spring Boot. La méthode \texttt{main()} appelle \texttt{SpringApplication.run()} qui :
\begin{enumerate}
  \item Scanne tous les packages sous \texttt{ma.fsb.sgrh}.
  \item Configure automatiquement les beans (auto-configuration).
  \item Démarre le serveur Tomcat intégré sur le port 8080.
\end{enumerate}

\section{Entités JPA (\texttt{entity/})}

Les entités sont des classes Java annotées \texttt{@Entity} qui correspondent directement aux tables MySQL. Lombok (\texttt{@Data}, \texttt{@Builder}) génère automatiquement les getters, setters, constructeurs et \texttt{toString()}.

\subsection{\texttt{Utilisateur.java}}
\begin{itemize}
  \item Implémente \texttt{UserDetails} (interface Spring Security) pour être directement utilisable dans l'authentification.
  \item Champs : \texttt{id}, \texttt{email} (unique), \texttt{motDePasse}, \texttt{role} (enum), \texttt{actif}.
  \item La méthode \texttt{getAuthorities()} retourne le rôle avec le préfixe \texttt{ROLE\_} requis par Spring Security.
  \item \texttt{isAccountNonLocked()} retourne la valeur de \texttt{actif}, permettant de bloquer un compte.
  \item Callbacks \texttt{@PrePersist} et \texttt{@PreUpdate} pour horodater automatiquement.
\end{itemize}

\subsection{\texttt{Employe.java}}
\begin{itemize}
  \item Hérite conceptuellement de \texttt{Utilisateur} via une relation \texttt{@OneToOne}.
  \item Contient les relations \texttt{@OneToMany} vers \texttt{Conge}, \texttt{Presence}, \texttt{FichePaie}, \texttt{Document}, \texttt{Anomalie} (cascade \texttt{ALL}).
  \item Champs métier : \texttt{matricule}, \texttt{nom}, \texttt{prenom}, \texttt{poste}, \texttt{departement}, \texttt{salaire}, \texttt{soldeConge}, \texttt{dateEmbauche}.
\end{itemize}

\subsection{\texttt{Conge.java}}
Demande de congé avec \texttt{dateDebut}, \texttt{dateFin}, \texttt{nbJours}, \texttt{typeConge} (enum), \texttt{statut} (enum, défaut \texttt{EN\_ATTENTE}), et \texttt{motif}.

\subsection{\texttt{Presence.java}}
Pointage quotidien avec \texttt{datePresence}, \texttt{heureArrivee} (\texttt{LocalTime}), \texttt{heureDepart}, et \texttt{estEnRetard} (booléen).

\subsection{\texttt{FichePaie.java}}
Fiche de paie mensuelle avec \texttt{mois}, \texttt{annee}, \texttt{salaireBase} et \texttt{salaireNet}.

\subsection{\texttt{Document.java}}
Document uploadé avec \texttt{nomFichier}, \texttt{typeDocument}, \texttt{cheminFichier}, \texttt{texteExtrait} (résultat OCR), et \texttt{statut} (défaut \texttt{EN\_COURS}).

\subsection{\texttt{Anomalie.java}}
Résultat d'analyse IA avec \texttt{scoreIsolation}, \texttt{typeAnomalie}, \texttt{niveauRisque}, \texttt{clusterLabel}, \texttt{dateDetection}, et \texttt{descriptionIa}.

\subsection{Enums}
Trois énumérations Java :
\begin{itemize}
  \item \texttt{Role} : \texttt{ADMIN}, \texttt{RH}, \texttt{EMPLOYE}.
  \item \texttt{StatutConge} : \texttt{EN\_ATTENTE}, \texttt{VALIDE}, \texttt{REFUSE}.
  \item \texttt{TypeConge} : \texttt{ANNUEL}, \texttt{MALADIE}, \texttt{SANS\_SOLDE}.
\end{itemize}

\section{Repositories (\texttt{repository/})}

Les repositories sont des interfaces qui étendent \texttt{JpaRepository<Entity, Long>}. Spring Data JPA génère automatiquement les implémentations SQL à partir des noms de méthodes.

\begin{center}
\begin{tabularx}{\textwidth}{|l|X|}
\hline
\textbf{Repository} & \textbf{Méthodes personnalisées} \\
\hline
\texttt{UtilisateurRepository} & \texttt{findByEmail()}, \texttt{existsByEmail()} \\
\hline
\texttt{EmployeRepository} & \texttt{findByMatricule()}, \texttt{findByUtilisateurId()}, \texttt{findByDepartement()} \\
\hline
\texttt{CongeRepository} & \texttt{findByEmployeId()}, \texttt{findByStatut()}, \texttt{findByEmployeIdAndStatut()} \\
\hline
\texttt{PresenceRepository} & \texttt{findByDatePresence()}, \texttt{findByEmployeIdAndDatePresenceBetween()}, \texttt{countByDatePresenceAndEstEnRetardTrue()} \\
\hline
\texttt{FichePaieRepository} & \texttt{findByEmployeId()}, \texttt{findByMoisAndAnnee()} \\
\hline
\texttt{DocumentRepository} & \texttt{findByEmployeId()}, \texttt{findByStatut()} \\
\hline
\texttt{AnomalieRepository} & \texttt{findByEmployeId()}, \texttt{findByNiveauRisque()}, \texttt{findTop10ByOrderByDateDetectionDesc()} \\
\hline
\end{tabularx}
\end{center}

\section{Sécurité (\texttt{security/})}

\subsection{\texttt{JwtService.java} --- Service JWT}
Gère le cycle de vie des tokens JWT :
\begin{itemize}
  \item \textbf{\texttt{generateToken()}} : crée un token signé HS256 contenant l'email comme \texttt{subject}, le rôle dans les claims, et une expiration de 24h.
  \item \textbf{\texttt{extractUsername()}} : décode le token et extrait l'email.
  \item \textbf{\texttt{isTokenValid()}} : vérifie que le token n'est pas expiré et que le username correspond.
  \item La clé de signature est décodée depuis Base64 et utilise HMAC-SHA256.
\end{itemize}

\subsection{\texttt{JwtAuthFilter.java} --- Filtre JWT}
Filtre Spring qui intercepte \textbf{chaque requête HTTP} :
\begin{enumerate}
  \item Extrait le header \texttt{Authorization}.
  \item Si le header commence par \texttt{Bearer }, extrait le token.
  \item Décode l'email depuis le token via \texttt{JwtService}.
  \item Charge l'utilisateur depuis la base de données.
  \item Vérifie la validité du token.
  \item Si valide, crée un objet \texttt{Authentication} et le place dans le \texttt{SecurityContext}.
\end{enumerate}

\subsection{\texttt{SecurityConfig.java} --- Configuration de sécurité}
Classe de configuration Spring Security :
\begin{itemize}
  \item \textbf{CORS} : autorise les requêtes depuis \texttt{http://localhost:3000} (frontend React).
  \item \textbf{CSRF} : désactivé (API stateless avec JWT).
  \item \textbf{Règles d'autorisation} :
  \begin{itemize}
    \item \texttt{/api/auth/**} : accès public (login, register).
    \item \texttt{/api/employes/**}, \texttt{/api/fiches-paie/**}, \texttt{/api/ia/**} : rôles ADMIN ou RH.
    \item \texttt{/api/utilisateurs/**} : ADMIN uniquement.
    \item Toutes les autres routes : authentifié.
  \end{itemize}
  \item \textbf{Session} : mode \texttt{STATELESS} (pas de session HTTP, authentification par token uniquement).
  \item \textbf{Password encoder} : \texttt{BCryptPasswordEncoder} pour le hashage des mots de passe.
\end{itemize}

\section{Contrôleurs (\texttt{controller/})}

\subsection{\texttt{AuthController.java}}
\begin{itemize}
  \item \textbf{POST \texttt{/api/auth/login}} : authentifie l'utilisateur, génère un token JWT, retourne le token et les infos utilisateur.
  \item \textbf{POST \texttt{/api/auth/register}} : crée un nouveau compte avec le rôle EMPLOYE par défaut.
\end{itemize}

\subsection{\texttt{EmployeController.java}}
CRUD complet sur les employés :
\begin{itemize}
  \item \textbf{GET} \texttt{/api/employes} : liste tous les employés.
  \item \textbf{GET} \texttt{/api/employes/\{id\}} : détail d'un employé.
  \item \textbf{POST} \texttt{/api/employes} : création (ADMIN/RH).
  \item \textbf{PUT} \texttt{/api/employes/\{id\}} : modification (ADMIN/RH).
  \item \textbf{DELETE} \texttt{/api/employes/\{id\}} : suppression (ADMIN uniquement).
\end{itemize}

\subsection{\texttt{CongeController.java}}
\begin{itemize}
  \item \textbf{GET} \texttt{/api/conges} et \texttt{/api/conges/employe/\{id\}} : liste des congés.
  \item \textbf{GET} \texttt{/api/conges/en-attente} : congés en attente de validation.
  \item \textbf{POST} \texttt{/api/conges} : soumission d'une demande (statut \texttt{EN\_ATTENTE}).
  \item \textbf{PUT} \texttt{/api/conges/\{id\}/valider} : valide le congé et \textbf{décrémente le solde} de l'employé.
  \item \textbf{PUT} \texttt{/api/conges/\{id\}/refuser} : refuse le congé.
\end{itemize}

\subsection{\texttt{PresenceController.java}}
\begin{itemize}
  \item \textbf{GET} par employé ou par date.
  \item \textbf{POST} \texttt{/api/presences} : enregistrement d'un pointage. Détecte automatiquement si l'employé est en retard (arrivée après 08h30).
\end{itemize}

\subsection{\texttt{DocumentController.java}}
\begin{itemize}
  \item \textbf{GET} \texttt{/api/documents} et \texttt{/api/documents/\{id\}}.
  \item \textbf{POST} \texttt{/api/documents/upload} : reçoit un fichier \texttt{multipart/form-data}, le sauvegarde dans le dossier \texttt{uploads/}, crée l'entrée en base avec statut \texttt{EN\_COURS}. L'appel à Tesseract OCR est prévu en asynchrone.
\end{itemize}

\subsection{\texttt{IAController.java}}
Pont entre le backend et le microservice IA Python :
\begin{itemize}
  \item \textbf{POST} \texttt{/api/ia/detect} : appelle \texttt{POST http://localhost:8000/detect} via WebClient.
  \item \textbf{POST} \texttt{/api/ia/cluster} : appelle \texttt{POST http://localhost:8000/cluster}.
  \item \textbf{GET} \texttt{/api/ia/rapport} : récupère le rapport consolidé depuis FastAPI.
  \item \textbf{GET} \texttt{/api/ia/anomalies} : retourne les 10 dernières anomalies depuis MySQL.
\end{itemize}

\subsection{\texttt{DashboardController.java}}
\begin{itemize}
  \item \textbf{GET} \texttt{/api/dashboard/stats} : retourne les 4 indicateurs clés (total employés, congés en attente, présents aujourd'hui, anomalies).
\end{itemize}

\section{DTOs (\texttt{dto/})}

\begin{itemize}
  \item \textbf{\texttt{AuthRequest}} : objet de requête login/register avec validation (\texttt{@Email}, \texttt{@NotBlank}).
  \item \textbf{\texttt{AuthResponse}} : objet de réponse contenant le token JWT et un sous-objet \texttt{UserDto} (id, email, role).
\end{itemize}


% ══════════════════════════════════════════════════════════════
\chapter{Microservice IA (\texttt{microservice-ia/})}
% ══════════════════════════════════════════════════════════════

\section{\texttt{requirements.txt}}
Dépendances Python du microservice :
\begin{center}
\begin{tabularx}{\textwidth}{|l|l|X|}
\hline
\textbf{Package} & \textbf{Version} & \textbf{Rôle} \\
\hline
\texttt{fastapi} & 0.104.1 & Framework web asynchrone haute performance \\
\hline
\texttt{uvicorn} & 0.24.0 & Serveur ASGI pour FastAPI \\
\hline
\texttt{scikit-learn} & 1.3.2 & Algorithmes ML (Isolation Forest, K-Means, StandardScaler) \\
\hline
\texttt{pandas} & 2.1.4 & Manipulation de données tabulaires (DataFrames) \\
\hline
\texttt{numpy} & 1.26.2 & Calcul numérique \\
\hline
\texttt{matplotlib} & 3.8.2 & Génération de graphiques pour les rapports \\
\hline
\texttt{pydantic} & 2.5.2 & Validation et sérialisation des données (modèles de requête/réponse) \\
\hline
\end{tabularx}
\end{center}

\section{\texttt{main.py} --- Serveur FastAPI}

\subsection{Configuration}
\begin{itemize}
  \item CORS activé pour \texttt{localhost:3000} (React) et \texttt{localhost:8080} (Spring Boot).
  \item Le serveur écoute sur le \textbf{port 8000}.
  \item Documentation Swagger automatique disponible sur \texttt{http://localhost:8000/docs}.
\end{itemize}

\subsection{Modèles Pydantic}
Modèles de validation des données entrantes et sortantes :
\begin{itemize}
  \item \textbf{\texttt{PresenceRecord}} : un enregistrement de présence avec \texttt{employe\_id}, \texttt{heure\_arrivee}, \texttt{heure\_depart}, \texttt{duree\_travail}, \texttt{est\_retard}, \texttt{nb\_absences\_semaine}.
  \item \textbf{\texttt{DetectRequest}} : liste de \texttt{PresenceRecord} + taux de contamination (défaut 5\%).
  \item \textbf{\texttt{ClusterRequest}} : liste de \texttt{PresenceRecord} + nombre de clusters (défaut 3).
  \item \textbf{\texttt{AnomalyResult}} : résultat pour un employé (score, anomalie, cluster, risque).
\end{itemize}

\subsection{Endpoint \texttt{POST /detect} --- Isolation Forest}
Algorithme de détection d'anomalies non supervisé :
\begin{enumerate}
  \item Reçoit un dataset JSON de présences (minimum 10 enregistrements).
  \item Extrait les 5 features : heure d'arrivée, heure de départ, durée de travail, retard, absences.
  \item Entraîne un modèle \texttt{IsolationForest} avec 100 estimateurs et un taux de contamination de 5\%.
  \item Calcule le \texttt{decision\_function} (score brut) et normalise entre 0 et 1 (1 = plus anormal).
  \item Applique en parallèle \texttt{K-Means} (k=3) après normalisation \texttt{StandardScaler}.
  \item Attribue un niveau de risque : score $> 0.7$ = ÉLEVÉ, $> 0.5$ = MOYEN, sinon FAIBLE.
  \item Retourne uniquement les anomalies détectées avec leur score, cluster et niveau de risque.
\end{enumerate}

\subsection{Endpoint \texttt{POST /cluster} --- K-Means}
Segmentation des profils comportementaux :
\begin{enumerate}
  \item Reçoit le dataset JSON.
  \item Normalise les features avec \texttt{StandardScaler} (moyenne 0, écart-type 1).
  \item Applique \texttt{K-Means} avec 3 clusters.
  \item Attribue un label à chaque cluster :
  \begin{itemize}
    \item \textbf{Cluster 0 - Normal} : présence régulière, peu de retards.
    \item \textbf{Cluster 1 - Retardataire} : retards fréquents, irrégularités d'horaires.
    \item \textbf{Cluster 2 - Absent} : taux d'absence élevé.
  \end{itemize}
  \item Retourne les labels par employé, les centres des clusters et la distribution.
\end{enumerate}

\subsection{Endpoint \texttt{GET /rapport}}
Retourne un rapport consolidé avec le résumé des analyses et des recommandations RH :
\begin{itemize}
  \item Nombre d'employés analysés et anomalies détectées.
  \item Distribution des clusters.
  \item Recommandations d'actions (entretiens, ajustement d'horaires).
\end{itemize}

\subsection{Endpoint \texttt{GET /generate-dataset}}
Génère un \textbf{dataset synthétique} pour les tests :
\begin{itemize}
  \item 100 employés $\times$ 180 jours = 18~000 enregistrements.
  \item 95\% des employés ont un comportement normal (arrivée $\sim$ 8h00, départ $\sim$ 17h00).
  \item 5\% (employés 1 à 5) ont un comportement anormal injecté (arrivée $\sim$ 9h30, départ $\sim$ 16h00, plus d'absences).
  \item Utilise \texttt{numpy.random} avec une seed fixe (42) pour la reproductibilité.
\end{itemize}


% ══════════════════════════════════════════════════════════════
\chapter{Fonctionnement de l'application}
% ══════════════════════════════════════════════════════════════

\section{Flux d'authentification (JWT)}

\begin{enumerate}
  \item L'utilisateur saisit son email et mot de passe sur la page \texttt{LoginPage}.
  \item Le frontend envoie un \texttt{POST /api/auth/login} avec les identifiants.
  \item Le backend vérifie les identifiants via \texttt{AuthenticationManager} (comparaison BCrypt).
  \item Si valide, \texttt{JwtService.generateToken()} crée un token JWT signé HS256 contenant l'email et le rôle, valide 24h.
  \item Le token et les infos utilisateur sont retournés au frontend.
  \item Le frontend stocke le token dans \texttt{localStorage} via \texttt{AuthContext.login()}.
  \item Pour chaque requête suivante, l'intercepteur Axios ajoute \texttt{Authorization: Bearer <token>}.
  \item Côté backend, \texttt{JwtAuthFilter} intercepte chaque requête, décode le token, et place l'authentification dans le \texttt{SecurityContext}.
\end{enumerate}

\section{Flux de gestion des congés}

\begin{enumerate}
  \item L'employé ouvre la page \texttt{CongesPage} et clique sur « Nouvelle demande ».
  \item Il remplit le formulaire (dates, type, motif) et soumet.
  \item Le frontend envoie un \texttt{POST /api/conges} avec le statut \texttt{EN\_ATTENTE}.
  \item Le congé est sauvegardé en base de données.
  \item Le responsable RH voit les congés en attente (filtré par statut).
  \item Il clique sur « Valider » $\rightarrow$ \texttt{PUT /api/conges/\{id\}/valider} :
  \begin{itemize}
    \item Le statut passe à \texttt{VALIDE}.
    \item Le \texttt{soldeConge} de l'employé est décrémenté de \texttt{nbJours}.
  \end{itemize}
  \item Ou il clique sur « Refuser » $\rightarrow$ le statut passe à \texttt{REFUSE}.
\end{enumerate}

\section{Flux de détection d'anomalies IA}

\begin{enumerate}
  \item Le responsable RH ouvre la page \texttt{IAPage} et clique sur « Lancer l'analyse ».
  \item Le frontend envoie un \texttt{POST /api/ia/detect} au backend Spring Boot.
  \item Le backend récupère les données de présence depuis MySQL.
  \item Le backend transmet ces données au microservice FastAPI via \texttt{WebClient} (\texttt{POST http://localhost:8000/detect}).
  \item FastAPI exécute l'algorithme \textbf{Isolation Forest} :
  \begin{itemize}
    \item Entraîne le modèle sur les 5 features de présence.
    \item Identifie les points aberrants (score d'anomalie normalisé).
  \end{itemize}
  \item En parallèle, FastAPI applique \textbf{K-Means} (k=3) pour segmenter les profils.
  \item Les résultats (anomalies + clusters) sont retournés au backend.
  \item Le backend sauvegarde les anomalies dans la table \texttt{anomalie} de MySQL.
  \item Les résultats sont envoyés au frontend qui les affiche :
  \begin{itemize}
    \item \textbf{Scatter plot} des clusters.
    \item \textbf{Tableau} des anomalies avec barres de score visuelles.
    \item \textbf{Anneau} de distribution des profils.
  \end{itemize}
\end{enumerate}

\section{Flux d'analyse OCR de documents}

\begin{enumerate}
  \item L'utilisateur ouvre \texttt{DocumentsPage} et clique sur « Uploader un document ».
  \item Il glisse un fichier (PDF, PNG ou JPG) dans la zone de drop ou le sélectionne.
  \item Le frontend envoie un \texttt{POST /api/documents/upload} en \texttt{multipart/form-data}.
  \item Le backend :
  \begin{enumerate}
    \item Sauvegarde le fichier physiquement dans le dossier \texttt{uploads/}.
    \item Crée une entrée en base avec le statut \texttt{EN\_COURS}.
    \item Invoque Tesseract OCR sur le fichier (en asynchrone).
  \end{enumerate}
  \item Tesseract analyse l'image/PDF et extrait le texte (français et arabe supportés).
  \item Le texte extrait est sauvegardé dans le champ \texttt{texte\_extrait} de la table \texttt{document}.
  \item Le statut passe à \texttt{TRAITE} (ou \texttt{ERREUR} si échec).
  \item L'utilisateur peut consulter le texte extrait dans le modal de détail.
\end{enumerate}

\section{Mode démonstration}

L'application intègre un \textbf{mode démo} qui permet de la présenter \textbf{sans aucun backend} :
\begin{itemize}
  \item Sur la page de login, trois boutons (Admin / RH / Employé) créent un utilisateur fictif avec un faux token.
  \item Toutes les pages contiennent des \textbf{données mock} (tableaux et graphiques pré-remplis).
  \item Les actions (ajout, validation, lancement IA) affichent des \textbf{notifications toast} de succès.
  \item Ce mode est idéal pour la \textbf{soutenance} si le backend ou la base de données ne sont pas disponibles.
\end{itemize}


% ══════════════════════════════════════════════════════════════
\chapter{Guide de démarrage}
% ══════════════════════════════════════════════════════════════

\section{Prérequis}

\begin{center}
\begin{tabularx}{\textwidth}{|l|l|X|}
\hline
\textbf{Outil} & \textbf{Version} & \textbf{Installation} \\
\hline
Node.js & $\geq$ 18.x & \url{https://nodejs.org} \\
\hline
Java JDK & 21 & \url{https://adoptium.net} \\
\hline
Python & 3.11+ & \url{https://python.org} \\
\hline
MySQL & 8.x & \url{https://dev.mysql.com/downloads/} \\
\hline
Tesseract & 5.x & \url{https://github.com/tesseract-ocr/tesseract} \\
\hline
\end{tabularx}
\end{center}

\section{Étapes de lancement}

\subsection{Étape 1 : Base de données}
\begin{lstlisting}[language=bash,title={Initialisation MySQL}]
mysql -u root -p < database/schema.sql
mysql -u root -p < database/seed.sql
\end{lstlisting}

\subsection{Étape 2 : Backend Spring Boot}
\begin{lstlisting}[language=bash,title={Lancement du backend (port 8080)}]
cd backend
./mvnw spring-boot:run
\end{lstlisting}

\subsection{Étape 3 : Microservice IA}
\begin{lstlisting}[language=bash,title={Lancement du microservice IA (port 8000)}]
cd microservice-ia
pip install -r requirements.txt
python main.py
\end{lstlisting}
La documentation Swagger est disponible sur \texttt{http://localhost:8000/docs}.

\subsection{Étape 4 : Frontend React}
\begin{lstlisting}[language=bash,title={Lancement du frontend (port 3000)}]
cd frontend
npm install
npm run dev
\end{lstlisting}

\section{Comptes de test}

\begin{center}
\begin{tabularx}{\textwidth}{|l|l|l|X|}
\hline
\textbf{Email} & \textbf{Mot de passe} & \textbf{Rôle} & \textbf{Accès} \\
\hline
\texttt{admin@sgrh.ma} & \texttt{password123} & ADMIN & Accès complet à toutes les fonctionnalités \\
\hline
\texttt{rh@sgrh.ma} & \texttt{password123} & RH & Gestion RH, IA, paie (pas la gestion des comptes) \\
\hline
\texttt{ahmed.benali@sgrh.ma} & \texttt{password123} & EMPLOYE & Congés, présences, documents personnels \\
\hline
\end{tabularx}
\end{center}


\end{document}
